\chapter{Throat cancer}

\section*{Overview}
Throat cancer is a condition in which abnormal cells grow and divide in an uncontrolled way. These cells can form a tumour, invade nearby tissues, or spread to other parts of the body through the blood and lymph systems. The behaviour, treatment, and outlook differ widely depending on the tissue of origin, the type of cells involved, and how far the disease has spread at the time of diagnosis.

\section*{Symptoms}
\begin{itemize}
    \item A new lump or thickening in the breast, testicle, skin, or elsewhere
    \item A change in the size, shape, colour, or feel of an existing mole or sore
    \item A sore that does not heal
    \item Persistent cough, hoarseness, or difficulty swallowing
    \item Unexplained weight loss, fevers, or night sweats
    \item Persistent tiredness and loss of appetite
    \item Unusual bleeding or discharge from any body opening
    \item Changes in bowel or bladder habits that persist
\end{itemize}

\section*{Causes}
Cancer arises from changes (mutations) in the DNA of cells. These changes may be inherited or acquired over a lifetime. Contributing factors include:

\begin{itemize}
    \item Tobacco smoking and exposure to second-hand smoke
    \item Excessive alcohol consumption
    \item Ultraviolet radiation from the sun or solariums
    \item Infections such as human papillomavirus, hepatitis B and C, and \textit{Helicobacter pylori}
    \item An unhealthy diet, physical inactivity, and obesity
    \item Exposure to cancer-causing chemicals such as asbestos, benzene, or certain workplace dusts
    \item Increasing age and a family history of cancer
\end{itemize}

\section*{Diagnosis}
\begin{itemize}
    \item Detailed medical history and physical examination
    \item Blood tests and tumour markers where appropriate
    \item Imaging such as X-ray, ultrasound, CT, MRI, or PET scans
    \item Biopsy, in which a sample of tissue is removed for laboratory examination
    \item Endoscopy or colonoscopy to look inside hollow organs
    \item Staging tests to determine whether, and how far, the disease has spread
\end{itemize}

\section*{Treatment}
Treatment is planned by a multidisciplinary team and depends on the type, stage, location, and the person's overall health. Common options include:

\begin{itemize}
    \item Surgery to remove the tumour and a margin of healthy tissue
    \item Radiation therapy to destroy cancer cells or shrink a tumour
    \item Chemotherapy, which uses drugs to kill cancer cells throughout the body
    \item Targeted therapy and immunotherapy, which act on specific features of cancer cells
    \item Hormone therapy for cancers that are sensitive to hormones
    \item Supportive and palliative care to manage symptoms and maintain quality of life
\end{itemize}

\section*{Prevention}
\begin{itemize}
    \item Do not smoke, and avoid exposure to second-hand smoke
    \item Maintain a healthy weight and be physically active
    \item Eat a balanced diet rich in vegetables, fruit, and whole grains
    \item Limit or avoid alcohol
    \item Protect skin from the sun and avoid solariums
    \item Take part in recommended vaccinations and cancer screening programs
    \item Follow workplace health and safety guidance to avoid carcinogens
\end{itemize}

\section*{When to Seek Medical Care}
See a doctor promptly for any persistent, unexplained symptom such as a new lump, a changing mole, unexplained weight loss, or unusual bleeding. Take part in the recommended screening for your age and sex. \textbf{Call 000} for severe symptoms such as difficulty breathing, chest pain, or sudden severe abdominal pain.

Always consult a qualified health professional for diagnosis and treatment of cancer; no information here replaces individual medical advice.
